\chapter{The accounting language}
\label{ch:language}

\section{What the current language can express}

A company model needs to keep disclosures, possible histories and chosen
methods distinct as its accounts evolve. A reported balance constrains a
history; it does not reveal every contract or transaction that generated it.
A collection policy supplies a behavioural assumption; it does not establish
what a customer will pay. An accounting method turns an eligible event into
state changes and postings; a balanced result alone does not establish that
the method applies. The language makes these distinctions explicit so that
reports retain the meaning of the facts and assumptions behind them.

The current design has several checked, bounded language profiles. Together
they cover source-constrained histories, conditional company paths, selected
shared-book operations and complete method authoring. Their interfaces are
explicit, but they are not yet one complete whole-company grammar.
Section~\ref{sec:common-program-authoring} shows the present composition
boundary: a single program contains complete method text and registries,
selects an inspectable state-invariant definition, and retains full event
intent beside the inputs used in a specialized calculation. The successor
in Section~\ref{sec:invariant-source-authoring} also writes the current finite
state invariants as explicit source, using the same typed evaluators. Wider
rule algorithms remain open, and a selected native event contract retains its
own implementation and applicability limits.

To represent the historical record, company method selection first binds to
a source-constrained partial state
(Section~\ref{sec:company-source-state}). Statement observations have
distinct relation roles rather than one undifferentiated set of totals
(Section~\ref{sec:statement-context-relations}). Selected balances are tested
jointly under explicit rounding and identity hypotheses
(Section~\ref{sec:joint-source-measurements}), with opening, movement and
closing dates kept separate across statements
(Section~\ref{sec:temporal-statement-network}). These declarations can expose
inconsistency or missing information without inventing the underlying books.

Several histories can remain consistent with those observations.
Section~\ref{sec:conditional-history-families} represents finite conditional
families with symbolic openings and explicit measurement attachment.
Typed economic quantities and constraints at each event give those histories
economic content (Section~\ref{sec:typed-economic-history}); cash and
quantity observations attach to the same history
(Section~\ref{sec:typed-history-observation-attachment}).
Unknown monetary amounts remain linked to selected method prerequisites
without choosing a single witness
(Section~\ref{sec:conditional-method-families}).
The retained return example follows every branch through settlement and sale,
including funding and intermediate observations
(Section~\ref{sec:native-conditional-sequence}). Outcome queries preserve
which results are possible or necessary and their exact finite bounds
(Section~\ref{sec:family-outcome-queries}). Stored definitions and results
can be recovered across restart, while the archive distinguishes stored
evidence from verified execution
(Section~\ref{sec:simulator-evidence-archive}).

Explicit future assumptions then determine how a conditional company path
evolves. The earliest matched text frontend handles the declared debt fragment.
The funded-year extension adds identified claim allocations, credit terms
and dated events; its imported annual projection still supplies the other
operating and accounting rules (Section~\ref{sec:funded-year-authoring}).
The operating-input example instead derives purchases, sales, claims,
expenses and closing statements from a historical opening, authored
quantities, prices, asset cohorts and contract terms
(Section~\ref{sec:company-input-authoring}). It imports no forecast journals:
a price or service-date change has consequences in the calculation under the
example's stated method restrictions.

Identified operations can repeat and carry unpaid claims and asset cohorts
into later years. Northstar's three-period path includes revised assumptions
and an interrupted year (Section~\ref{sec:company-operation-paths}).
Its text frontend preserves the same typed records, named settlement
allocations and incomplete paths
(Section~\ref{sec:company-path-authoring}).
Named, unweighted scenarios keep the alternative assumptions separate.
Their collection policies have pure typed expressions and bounded reusable
functions with matched text declarations
(Section~\ref{sec:company-family-language}). Function calls preserve units,
company context and the distinction between declared and actually used
inputs. They reuse a calculation without supplying missing facts or creating
endogenous events.

Combining methods requires compatible economic state as well as account names.
Portable module identities expose the additional state required by a
shared-book return method (Section~\ref{sec:company-module-language});
its matched authoring and finite execution appear in
Section~\ref{sec:company-return-reference}. Heterogeneous lease rights require
explicit account allocations and funding conditions
(Section~\ref{sec:lease-shared-book}). Aggregate balances may leave those
allocations unidentified, producing different complete or interrupted
outcomes (Section~\ref{sec:lease-allocation-family}).
Actual rounded lease disclosures constrain some possibilities before
individual contract state is known
(Section~\ref{sec:henkel-lease-measurements}).

A common caller does not establish a common book. The return-and-lease example
first establishes compatible state, then follows both kinds of event through
one declared cash balance (Section~\ref{sec:shared-return-lease-book}).
Its recorded statements reconcile while retaining unresolved tax assessment
and incomplete fiscal coverage
(Section~\ref{sec:recorded-shared-statements}). The complete opening and
mixed event sequence have matched text and typed forms
(Section~\ref{sec:shared-book-authoring}).
Section~\ref{sec:shared-tax-information} explains why the original tax
method refuses that mixed state, then develops explicitly conditional
successors for assessment, payments, period carry and annual programs.
Those successors preserve unknown declarations and stopped prefixes. They
do not establish the company's actual tax policy or complete fiscal history.

Methods state their prerequisites before applying financial effects.
Reusable selected requirements have matched typed and text declarations
(Section~\ref{sec:method-requirement-language}).
Exact input binding, dimensional products and guarded applicability require
explicit version selection (Section~\ref{sec:requirement-successors}).
Collection coverage, identity order, signed totals and state-domain findings
add constraints that scalar types cannot express
(Section~\ref{sec:collection-state-requirements}).
Typed calculations and atomic state updates then support exact source-bound
replay of the retained return example
(Section~\ref{sec:typed-return-transitions}).
Complete method text includes these prerequisites, monetary calculations,
postings and memory changes
(Section~\ref{sec:complete-method-authoring}).
The common program preserves that source and its explicit owners; it does
not silently replace a newly written method with a named demonstration.

The historical AccDSL Core 0 proposal remains the teaching baseline. It
describes entities and books, accounts, observations, shared unknowns,
constraints, alternative interpretations, contract instances, accounting
rules and reporting mappings, with pure expressions and a restricted posting
effect. Appendix~\ref{app:contracts} preserves its grammar and examples.
The proposed portable records and first company profile refine that baseline.
Numerical and forecasting algorithms can live in typed plugins; unrestricted
numerical code need not be reimplemented in the DSL. General recursive or
simultaneous algorithms, endogenous decisions, source admission and broader
accounting profiles remain separate language and method work.

The historical Core 0 compiler has six ordered stages: parsing with source
locations; name and module resolution; type and temporal-context checking;
lowering to the IR; constraint and rule-applicability analysis; and
execution-plan construction. It never silently supplies an economic
assumption to make a program compile. A valid historical representation can
still require additional assumptions before a particular simulation can run.
The proposed portable compilation sequence is specified separately in
Section~\ref{sec:portable-authoring}. Structural compilation of a current
program likewise does not execute its financial methods or resolve unknown
eligibility findings.

The IR and evaluator define meaning separately from surface spelling.
Formatting that preserves the lowered method preserves its semantic digest.
The complete authored program also retains the original method source, so its
program digest changes when those bytes change. In the checked common-program
example, adding one space leaves the lowered method and materialized financial
request identical while changing the source and program digests. Source spans
and reviewer decisions can change evidence provenance without changing the
numerical model. Evidence, semantic-model and execution identities therefore
remain distinct; a field named \code{hash} must identify what it actually binds.

For the preserved Core 0 examples, the complete grammar is
\code{spec/core-0.lark}; the original design directory retains
\code{spec/core-ir.schema.json} and \code{spec/core-contract.md}.
They define reference resolution, expression arity, type and context checks,
operations, failures and identity. A structural schema check alone does not
establish those semantic properties. This revision uses ASCII identifiers and
UTF-8 strings, with explicit lexical rules for currency suffixes, dates and
decimal literals. Parsing the complete demonstration checks the grammar; it
does not establish an implemented compiler or solver. The following type
distinctions apply before any of those operations can be interpreted.

\section{Types that protect meaning}

The following table uses the preserved Core 0 type vocabulary. Its distinctions
motivate the more precise portable types introduced below.

\begin{longtable}{L{40mm}L{96mm}}
\toprule Type & Meaning and checks\\\midrule\endhead
\code{Money<C>} & Exact amount in currency C. Addition requires identical
currency; conversion requires a dated rate and policy.\\
\code{Quantity<U>} & Units of shares, goods or other commodities. Quantities
cannot be added to money or used as a functional-currency posting.\\
\code{Rate<B>} and \code{YearFraction} & Rate with basis B, plus a separately
computed accrual fraction. Simple and compound accrual functions are distinct.\\
\code{Instant}, \code{Period}, \code{Calendar} & Instant observations and
flows over intervals; interval operations preserve calendar conventions.\\
\code{Account<K,Book>} & Account role K and owning book. Cross-book effects
require separate balanced batches linked by an economic event.\\
\code{Observation<T>} & Source assertion with contextual and precision
metadata; accessing its value retains its constraint and provenance.\\
\code{Symbol<T,D>} & Shared unknown with domain D; repeated references denote
the same value in a history.\\
\code{Choice<E>} & A finite structural choice over enumeration E. No probability
is implied.\\
\code{Distribution<T>} & Explicit probability measure with parameter,
conditioning and measure metadata. It is not a Symbol or a Choice.\\
\code{SourceRef}, \code{RuleRef}, \code{ModelRef} & Versioned evidence or
implementation references; unresolved references prevent concrete execution.\\
\code{Measurement<T>} & Value plus valuation date, basis, method, assumptions
and numerical status. Recognition remains a rule decision.\\
\bottomrule
\end{longtable}

The portable profile spells the rate type \code{Rate<convention>} and makes
both calendar and accrual convention explicit in
\code{YearFraction<calendar, convention>}. \code{InstantValue<T>} and
\code{Flow<T, [start,end)>} distinguish dated stocks from interval flows;
their context carries the calendar. An approximate numerical component of a
historical \code{Measurement<T>} uses \code{Approximation} with its method,
error description and environment references. An exact measurement remains
exact; valuation context and recognition authority are separate in either case.

Core expressions use exact integers and rationals, Boolean logic, enumerations,
dates, records, lists, comparisons and registered pure functions. A decimal
literal is parsed exactly, not first converted through binary floating point.
Numerical plugins can return approximate measurements, but must declare their
error or numerical status and the rounding used when posting money.

Missingness is a tagged type with distinct values: \code{notDisclosed},
\code{notApplicable}, \code{unreadable}, \code{withheld}, and
\code{unresolvedMeaning}. Zero is an ordinary known number. A missing datum can
be replaced by a declared Symbol or assumption in a model branch; the source
observation stays missing. An unknown is not necessarily random, and an
estimate is not necessarily missing.
On portable import, \code{notDisclosed}, \code{unreadable}, \code{withheld} and
\code{unresolvedMeaning} become \code{Unknown} with the corresponding reason
and obligation references. \code{notApplicable} becomes the explicit
\code{NotApplicable} variant, never an unknown or zero.

Addition of monetary observations also requires compatible entity, perimeter,
time, dimensions and basis. These checks operate on contextual expressions in
addition to the currency type. An analyst can deliberately aggregate dimensions
or transform a basis using an explicit function. For example,
\code{sumOver(segment, ...)} records the dimension being removed and requires
the segment partition to be declared; it does not assume every disclosed
segment total is additive.

\input{sections/portable-values}
\section{A small model with a continuous interpretation family}

The following source uses a synthetic observation whose full context is stored
in the evidence registry. The example imports the demonstration debt package,
which uses one-year simple accrual, no fees and immediate rate resets. It is
not an IFRS rule package.

All dependencies are instantiated in
\code{fixtures/demo-debt-package.json}: the source reading, balanced opening
state, currency precision, contract state and dates, reset and interest
operations, closing and report mappings, request and expected outputs.
\code{fixtures/borrowing.acc} is the complete source, including the retained
earnings account and an \code{opening} declaration. The shorter listings below
explain its declarations in stages. They do not replace that full source.

\begin{lstlisting}
module BorrowingExample version "1";
use DemoDebt version "1";
source note7 = external("synthetic-note-7");
entity Acme functional USD;
book Local entity Acme basis DemoDebt version "1";
account Cash : Asset book Local;
account Debt : Liability book Local;
account InterestExpense : Expense book Local;

observation debt : Money<USD> instant @2025-12-31
  = 100 USD from note7;
symbol alpha : Real in [0, 1]
  because "The fixed-rate proportion is not disclosed";
\end{lstlisting}

The remaining declarations use this same unknown fraction in both loan
components. A change to the fraction therefore changes the whole debt mix
consistently.

\begin{lstlisting}
let fixedPrincipal = alpha * value(debt);
let floatingPrincipal = (1 - alpha) * value(debt);
constraint split:
  fixedPrincipal + floatingPrincipal == value(debt);

contract loans : DebtMix book Local {
  principal = value(debt);
  fixed_fraction = alpha;
  fixed_rate = 0.05 per year;
  initial_floating_rate = 0.05 per year;
}
alternative allFixed { assume alpha == 1; }
alternative allFloating { assume alpha == 0; }
scenario rateShock {
  set market_rate = constant(0.07 per year);
}
\end{lstlisting}

The parent model retains every $\alpha\in[0,1]$. The two named alternatives
are useful conditional views of that family, not a declaration that the
endpoints exhaust it. Requesting the parent family's exact bounds produces
$[5,7]$ for the next year's expense under \code{rateShock}. Requesting a
conditional run of \code{allFixed} supplies $\alpha=1$. Requesting a weighted
mixture requires a separate weight model; the two labels do not create equal
weights.

An observation declaration is a reference to evidence plus a proposed typed
reading. The compiler checks that the registry context agrees with the declared
period and type. If the evidence is rounded, \code{value(debt)} denotes a
shared constrained value rather than a silently exact 100. This synthetic
registry explicitly declares an exact value, which keeps the illustration
deterministic.

\begin{minipage}{\linewidth}
\section{Posting rules and explicit exceptions}

An event type declares the information needed by a rule. The example below
assumes an event representing interest accrued and paid on the same date.
Separate accrual and settlement events are needed when those dates differ.

\begin{lstlisting}
event InterestPaid(principal: Money<USD>,
                   rate: Rate<annual_simple>,
                   fraction: YearFraction) book Local;

rule recogniseInterest for InterestPaid as e
when e.principal >= 0 USD {
  let charge = roundMoney(simpleInterest(e.principal, e.rate, e.fraction),
                         currencyScale(USD), "half_even");
  require charge >= 0 USD;
  post {
    debit InterestExpense charge;
    credit Cash charge;
  }
}
\end{lstlisting}
\end{minipage}

The rule cannot perform file I/O, mutate another entity or sample a random
number. It returns a proposed batch. The runtime verifies balance and all
preconditions before committing state. Negative interest requires an explicit
supported rule; it is not mishandled by this intentionally restricted example.
Rounding is a declared demonstration convention, not a universal accounting
default.

Rules can declare \code{overrides ruleName}. The resulting override graph
must be acyclic. For a concrete event, the evaluator determines all applicable
rules and removes any that are overridden by another applicable rule. Exactly
one maximal rule must remain for a required recognition operation. An empty
set of applicable rules produces \code{UnsupportedTreatment}; more than one
maximal rule produces
\code{RuleConflict}. File order and a hidden numeric priority never break ties.

If a guard depends on a symbolic value, the solver partitions the current
feasible family into the guard and its negation, subject to the existing
constraints. This is a logical partition, not a random draw. An incomplete
solver returns \code{UnknownApplicability}, not a guessed rule. Statically
provable disjointness or total coverage can be cached; otherwise the runtime
performs the same checks on every concrete event.

\section{Alternative readings, assumptions and updates}

An alternative adds constraints and possibly a supported semantic mapping to
its parent model. Its record contains a rationale, supporting spans, opposing
evidence, author, creation time and scope. Child alternatives inherit all
parent constraints; they cannot quietly discard an inconvenient disclosure.
A deliberate source rejection is a separate review decision with an explicit
reason and a new evidence view.

Two independent alternatives may coexist, while two mutually exclusive readings
of the same text require an explicit choice group. The representation stores
these relationships in a constraint graph. It does not take a Cartesian
product blindly. For example, interpreting an instrument as a liability affects
both its balance-sheet account and the treatment of its periodic payments;
those choices must be linked through one shared treatment variable.

New evidence narrows an existing family by conjunction when it concerns the
same definitions and reporting basis. If it changes the basis or corrects an
earlier assertion, the system creates a revised evidence view and records the
relationship. Posterior reweighting, when enabled, is a separate operation
with likelihood provenance. A branch that becomes inconsistent is retained
as historically considered and rejected, with the evidence that rejected it.

\section{Modules and stable intermediate representation}

A package identifies its namespace, semantic version, dependencies, effective
dates, jurisdiction, supported entity classes, recognition operations and
measurement policies. Standard packages define allowable treatments; industry
packages add economic objects and specialised rules; company packages map
source concepts and declare evidenced elections. A company extension cannot
override a mandatory standard rule merely because its package is loaded last.
An override is valid only at a declared extension point and within the selected
standard's policy domain.

The IR has separate node classes for evidence, contextual observations,
expressions, domains, constraints, alternatives, contracts, economic events,
measurement results, rules, journal batches, report mappings and runs.
Each node has a stable identifier and explicit input references. No expression
stores executable host-language text. Plugins are referenced by versioned
identifiers with typed signatures. This permits a second implementation to
execute the same supported program without evaluating arbitrary embedded code.

The IR serialisation version is independent of the DSL version and of any
accounting-standard edition. Migrations preserve the original payload and
record changed meanings. A syntax migration can be automatic when semantics
are unchanged. A changed recognition policy requires a new model version,
even if the source file still parses.

\begin{minipage}{\linewidth}
\section{Errors are part of the interface}

The compiler and runtime return typed diagnostics with a source span,
affected entity and period, dependency chain and suggested local action.
The principal diagnostic categories have distinct recovery paths:

\begin{center}
\begin{tabular}{L{61mm}L{75mm}}
\toprule Diagnostic & Meaning and next action\\\midrule
\code{TypeMismatch}, \code{ContextMismatch} & Resolve incompatible units,
periods, entities or measurement bases; do not coerce implicitly.\\
\code{MissingEvidence} & Supply a source or label a model input as an assumption.\\
\code{InfeasibleHistory} & Inspect the conflicting evidence and model
constraints; preserve the rejected candidate.\\
\code{RuleConflict}, \code{UnknownApplicability} & Resolve rule scope or
preserve a justified conditional partition.\\
\code{UnsupportedTreatment} & Add a reviewed capability or leave the affected
output incomplete.\\
\code{UnbalancedBatch} & Reject the atomic update and diagnose the posting rule.\\
\code{SolverIncomplete}, \code{NumericalFailure} & Preserve the target and
inputs; repair or change the computational method.\\
\code{InsufficientSimulationState} & Add declared conditional assumptions or
request a family of results.\\\bottomrule
\end{tabular}
\end{center}

Warnings can describe low extraction confidence or an approximate economic
model, but they do not replace these errors. A required unknown input prevents
a concrete unconditional result. The analyst can add a declared conditional
assumption, request bounds, or leave that output incomplete. Every path has a
different, explicit result type; there is no generic fallback to zero.
\end{minipage}\par
